% ==========================================================================
% Models of cancer dynamics under therapy
% ==========================================================================

\section{Models of Cancer Dynamics Under Therapy}
\label{sec:models}
\subsection{Spatial Model}
There is a plethora of cancer models based on various modeling paradigms
(see e.g., \cite{cristini2010multiscale}, \cite{Sekyere2025VascularTumor}, \cite{wcislo20093}). Most of them concern only the
phase of tumor growth. Much less consider various anti-cancer therapies
(see e.g. \cite{Sujitha2023RadiotherapyBrainTumor}, \cite{dzwinel2021supermodeling}). Let us assume that we have an array of such
models ordered in decreasing complexity starting from the reality R: R
\(\rightarrow\)$M_{\infty} \to \ldots \to M_i \to M_{i-1} \to \ldots \to M_1 \to M_0$.
We assume that the number of parameters will decrease from the most
complicated to the simplest model. Let us assume that $M_{\infty}$
is the most faithful metaphor of reality R, while possible differences
come from various sources of uncertainty. In general, even in reality,
assuming the same two initial setups, one can obtain (after some time)
different outcomes resulting from the nonlinear amplification of small
variations in the systems (chaotic behavior). Due to the measurement,
numerical, and round-off errors the same situation may happen for
computer simulations. So, here, before introducing the reduced formulation, we consider a more
general radiotherapy-driven reaction-diffusion model $M_{2}$ in which tumor
evolution is coupled with treatment response, oxygen-dependent
radiosensitivity, and unresolved microenvironmental effects. In its
compact form, such a model may be written as \cite{Sujitha2023RadiotherapyBrainTumor}:
\begin{equation}
\label{eq:pde_general}
\begin{split}
    \frac{\partial u(x,t)}{\partial t}
&=
\nabla \cdot \left( D(x)\nabla u \right)
+
F_{b}(u,x,t;\theta_{b})
-
U(t)\,G_{R}(u,R(x,t),h(x,t);\theta_{R}) \\
&+
\mathcal{M}(u,x,t).
\end{split}
\end{equation}
where \(u(x,t)\) denotes the local tumor cell density, \(D(x)\) is a
possibly heterogeneous diffusion coefficient, \(F\) represents intrinsic
tumor growth and tissue-limited proliferation, \(G_{R}\) describes
radiotherapy-induced cell killing modulated by the delivered dose
\(R(x,t)\)and local radiosensitivity factors \(h(x,t)\), while
\(\mathcal{M}\) collects additional microenvironmental mechanisms such as
hypoxia, vascular remodeling, delayed cell death, repair processes, and
repopulation between treatment fractions. Additionally, \(U(t)\) defines
the radiotherapy fractionation schedule, i.e. the sequence of
irradiation events. A convenient representation is
\begin{equation}
\label{eq:therapy_schedule}
U(t)
=
\sum_{k=1}^{N_f}
\delta_{\tau}(t-t_k).
\end{equation}
where \(N_{f}\) is the total number of fractions, \(t_{k}\) denotes the
time of the \(k\)-th irradiation, and \(\delta_{\tau}(t - t_{k})\) is a
narrow temporal pulse representing dose delivery over a short time
interval. A simplified instance of this general formulation is then obtained by
assuming logistic tumor growth, constant diffusion, multiplicative
radiotherapy-induced killing, and representing all unresolved effects by
a residual term \cite{Sujitha2024ChemoRadiotherapy}:
\begin{equation}
\label{eq:pde_simplified}
\frac{\partial u(x,t)}{\partial t}
=
D\nabla^{2}u
+
\rho u\left(1-\frac{u}{K}\right)
-
\beta U(t)R(x)H(x)u
+
\xi(u,x,t).
\end{equation}
where: \(\rho\) denotes the intrinsic tumor proliferation rate,
\(K\) is the effective carrying capacity,
\(\beta\) Radiosensitivity coefficient, the residual term \(\xi(u,x,t)\) represents unresolved biological
mechanisms and modeling uncertainties.

\subsection{ODE Model}

In practice, 3D tumor structure can be observed using imaging techniques
such as CT or MRI; however, these observations are limited to discrete
snapshots separated by relatively large time intervals. Consequently,
the full spatio-temporal field \(u(x,t)\) is not directly observable.
Moreover, clinical measurements typically capture only the visible tumor
volume, which may underestimate the true tumor burden due to limited
resolution and the presence of diffuse or microscopic disease. In
reduced models, this discrepancy can be partially compensated by
introducing an effective tumor mass, for example by incorporating safety
margins or scaling factors. So, instead, measurements typically
correspond to macroscopic quantities. Therefore, we introduce an
observable such as tumor mass as:
\begin{equation}
\label{eq:mass}
y(t)
=
\int_{\Omega} u(x,t)\,dx
\qquad \text{and} \qquad
\int_{\Omega} \nabla^{2}u\,dx = 0
\end{equation}
assuming zero-flux BC. Combining Equation (\ref{eq:pde_simplified}) and  Equation (\ref{eq:mass}):
\begin{equation}
\dot y=
\rho y
-\frac{\rho}{K}\int_\Omega u^2dx
-\beta\int_\Omega URHu\,dx
+\int_\Omega \xi\,dx .
\end{equation}
If we now wish to approximate it using a logistic ODE:
\(ay - by^{2} = ay(1 - \frac{b}{a}y)\) in combination with radiotherapy,
e.g.
\begin{equation}
\small
\begin{aligned}
\dot y
&=
\dot y = ay - by^2 - cU(t)\bar R_{\mathrm{eff}}(t)y
\\
\delta(t)
&=
-\frac{\rho}{K}
\Bigl(\int_\Omega u^2dx-\alpha y^2\Bigr)
-\beta
\Bigl(
\int_\Omega U(t)RHu\,dx
-\bar R_{\mathrm{eff}}(t)y
\Bigr)
+\int_\Omega \xi(u,x,t)\,dx .
\end{aligned}
\end{equation}
where:
\begin{itemize}
    \item $R(x,t)$ local radiotherapy at PDE.

    \item $\bar R_{\mathrm{eff}}(t) = \frac{\int_{\Omega}^{}{R(x,t)H(x)u(x,t)\text{\,}dx}}{\int_{\Omega}^{}{u(x,t)\text{\,}dx}}$ - state-dependent exact effective mean

    \item the available macroscopic input used in ODE/NODE.
    
\end{itemize}

This is a very powerful example, as it highlights three sources of
error:
\begin{enumerate}
\item \emph{Spatial non-linearity}
\begin{equation}
\int_{\Omega} u^{2}\,dx
\neq
\alpha y^{2}
\end{equation}
\item \emph{Inconsistent response to radiotherapy}
\begin{equation}
\int_{\Omega} R(x,t)H(x)u\,dx
\neq
\bar{R}_{\mathrm{eff}}(t)\,y
\end{equation}
\item \emph{Missing physics / biology}
\begin{equation}
\int_{\Omega} \xi(u,x,t)\,dx
\neq 0
\end{equation}
\end{enumerate}
Let the ODE surrogate therefore take the ODE system form:

\begin{equation}
\label{eq:pinode_residual}
\frac{dy}{dt}
=
f_{RT}\!\left(y,U(t);\theta\right)
+
g_{\phi}\!\left(y,t,U(t)\right)
\end{equation}
\(y(t)\) - the system state (e.g., tumor mass, or an extended
state in NODE),
\(f_{RT}\) - the known mechanistic core (radiotherapy-driven dynamics),
\(g_{\phi}\) - a learned correction term accounting for unresolved
effects,
\(U(t)\) - the control/input function, describing the applied therapy
over time.

In our case:
\begin{equation}
\label{eq:frt}
f_{RT}(y,U(t);\theta)
=
\rho\,y(1-y)
-
\beta\,\overline{R}_{\mathrm{eff}}(t)\,U(t)\,y
\end{equation}
The $g_{\phi}$ -- learned correction, which depends on: the
complete cell distribution, the complete dose distribution, the
microenvironment, treatment history. Therefore, it compensates for:
(a) spatial heterogeneity of the dose,
(b) hypoxia,
(c) delayed cell death,
(d) repopulation between fractions,
and (e) other unmodelled processes.
Thus, \(g_{\phi}\) accounts for unresolved effects, while the treatment
itself is represented explicitly by the control function \(U(t)\).
